\section{Introduzione}
Lo scopo di questo elaborato è quello di analizzare come le tecniche offensive e difensive nelle console di gioco casalinghe si siano evolute nel tempo.\\
Il modello di minaccia che caratterizza questi sistemi è il modello MATE -- \textit{Man At The End}, ovvero quello in cui l'attaccante ha completa disposizione fisica della macchina.\\
Analizziamo quindi lo stato dell'arte delle tecniche di difesa impiegati dalle più grandi case produttrici di console di gioco, e di come siano comunque state prese di mira attraverso exploit software, corruzione della memoria e attacchi hardware come la \textit{Fault Injection}, con l'obiettivo di aggirare la \textit{Chain of Trust}.

\subsection{Dispositivi Chiusi e l'Ecosistema \textit{Walled Garden}}

Il mercato dell'elettronica di consumo si divide storicamente in piattaforme aperte (PC Windows/Linux) e piattaforme chiuse (console, Set-Top Box, dispositivi iOS). Una tendenza consolidata mostra che l'apertura delle piattaforme attira un numero maggiore di sviluppatori e consumatori, ma abbassa drasticamente la soglia di competenza richiesta a un attaccante per comprometterle.

Per mantenere il controllo e proteggere la proprietà intellettuale, il cui valore commerciale risiede principalmente nel software e non nell'hardware,ci produttori adottano il paradigma del Dispositivo Chiuso, creando un ecosistema iper-controllato noto come Walled Garden (Giardino Murato). In questo ambiente, il produttore impone vincoli crittografici stringenti affinché l'hardware accetti esclusivamente codice firmato e autorizzato, negando all'utente finale qualsiasi privilegio di amministrazione.

\subsection{Dal Modello Black-Box al Modello White-Box}

La sicurezza informatica classica si fonda su assunzioni crittografiche e attacchi di tipo Black-Box. In tale scenario, si assume che i due attori legittimi (Alice e Bob) abbiano il controllo esclusivo dei propri dispositivi: la minaccia principale è un agente esterno che intercetta i dati in transito (attacco Man-In-The-Middle).

Nelle console questa assunzione crolla: ci troviamo di fronte a un modello White-Box in cui l'attaccante coincide con l'utente stesso. Bob possiede fisicamente la macchina e il software di Alice, ribaltando completamente il paradigma difensivo.

\subsection{L'Attacco MATE (Man-At-The-End)}

Questa specifica minaccia è definita nella letteratura accademica come attacco Man-At-The-End (MATE). Il paper fondativo sul tema è \hyperref[fonte:falcarin]{\textbf{Falcarin, Collberg et al. (2011)}}~[\ref{fonte:falcarin}]. In uno scenario MATE, il dispositivo hardware e l'ambiente di esecuzione non sono considerati fidati (Untrusted): l'attaccante ha accesso diretto alla macchina e al software in qualsiasi condizione, a prescindere dal fatto che il software sia in esecuzione o meno.

\begin{quote}
\textit{\small "In a MATE scenario, the attacker has unlimited powers: they can trace every instruction, inspect all memory and cache, halt execution at any moment, alter code or memory at will, and do all this for as long as desired, in collusion with as many other attackers as they can find."} --- Falcarin, Collberg et al. --- Software Protection: A look at MATE attacks, 2011~[\ref{fonte:falcarin}]
\end{quote}

A causa di questa sproporzione di poteri, l'attaccante MATE gode di capacità che superano di gran lunga quelle del difensore. Un attaccante MATE può in particolare:

\begin{enumerate}
\item Tracciare ogni singola istruzione del programma in esecuzione.
\item Visualizzare e modificare il contenuto della memoria, dei registri e della cache della CPU.
\item Interrompere l'esecuzione in qualsiasi momento e analizzarla in modo offline (tramite dump di memoria).
\item Alterare il codice sorgente o i binari, e collaborare con altri attaccanti per condividere le scoperte.
\end{enumerate}

\subsection{La Chain of Trust e la Root of Trust}

Poiché il software da solo non può difendersi in un ambiente ostile in cui l'utente può alterare la memoria a piacimento, i Walled Garden si affidano alla sicurezza hardware per fondare una Trusted Computing Base. Come documentato nel survey di \hyperref[fonte:edgar]{\textbf{Edgar e Manz (2017)}}~[\ref{fonte:edgar}], questa base prende la forma di una Chain of Trust (Catena di Fiducia) crittografica:

\begin{enumerate}
\item Root of Trust (BootROM): Il primissimo codice eseguito all'accensione è inciso in modo immutabile nel silicio del processore. Non può essere modificato, nemmeno dal produttore.
\item Caricamento a cascata verificato: La BootROM verifica la firma digitale del componente successivo (il Bootloader) prima di passargli il controllo. Il Bootloader verifica il Kernel, e così via fino al gioco.
\end{enumerate}

\begin{center}
BootROM  \(\to\)  Bootloader  \(\to\)  Kernel (S.O.)  \(\to\)  Gioco
\end{center}

Se in qualsiasi momento un anello della catena rileva che il codice successivo è stato alterato (Tampering), il processo di avvio viene interrotto. L'intero filone dell'hacking delle console si basa sul tentativo dell'attaccante MATE di corrompere anche uno solo di questi anelli.